\section{Implement a Min Heap}
\label{sec:heap-implement}

\problemstatement{Implement a min-heap supporting three operations:
$\textit{push}(x)$ (insert a new value), $\textit{pop}()$ (remove and
return the smallest value), and $\textit{top}()$ (peek at the smallest
value without removing it), each in $O(\log n)$ time for push/pop and
$O(1)$ for top.}

\begin{patternbox}
This section has no ``keyword'' in the usual sense -- it is the mechanical
foundation every other section in this chapter builds on. The skill being
tested is not pattern recognition but the array-indexing arithmetic that
lets a heap live inside a flat array without ever needing explicit
left/right/parent pointers.
\end{patternbox}

\subsection*{Understanding the array-as-tree representation}

Store the heap's elements in a plain array $\textit{heap}[0..n-1]$,
interpreted as a complete binary tree read level by level, left to right.
For a node at index $i$ (0-indexed):
\[
  \textit{parent}(i) = \left\lfloor \frac{i-1}{2} \right\rfloor, \qquad
  \textit{left}(i) = 2i+1, \qquad \textit{right}(i) = 2i+2.
\]
This works because a complete binary tree (every level full except
possibly the last, which fills left to right with no gaps) has exactly
this relationship between a node's position in level order and its
children's positions -- no pointers are needed, only arithmetic.

\begin{ideabox}[Sift Up (used by push)]
Append the new value at index $n$ (the next free slot). While this node's
value is smaller than its parent's value, swap them and continue checking
from the parent's old position. This restores the min-heap invariant
locally along the single path from the new leaf up to the root, which is
always sufficient: every other path in the tree was already valid before
the insertion and is untouched by these swaps.
\end{ideabox}

\begin{ideabox}[Sift Down (used by pop)]
Save the root's value (this is the return value). Move the \emph{last}
element into the root's slot and shrink the array by one. While this
relocated value is larger than the smaller of its two children, swap it
with that smaller child and continue checking from the new position. This
restores the invariant locally along a single path from the root downward.
\end{ideabox}

\subsection*{Why sifting along a single path is always sufficient}

Both operations rely on the same structural fact: the min-heap invariant
is a \emph{local} condition (each node versus its own parent/children),
and inserting or removing one element only disturbs the invariant along
\emph{one} root-to-leaf path -- the path the new or relocated element
travels. Every other node's relationship to its own parent and children is
completely unaffected, so checking and fixing violations one step at a
time along that single path, stopping the instant no violation remains, is
guaranteed to restore a fully valid heap. This is also exactly why each
operation costs $O(\log n)$: a path from root to leaf in a complete binary
tree with $n$ nodes has length $O(\log n)$, since the tree's height is
$\lfloor \log_2 n \rfloor$.

\subsection*{Algorithm}

\textbf{Push($x$):}
\begin{enumerate}
  \item Append $x$ at index $n$; let $i \gets n$; increment $n$.
  \item While $i > 0$ and $\textit{heap}[\textit{parent}(i)] >
        \textit{heap}[i]$: swap them; $i \gets \textit{parent}(i)$.
\end{enumerate}

\textbf{Pop():}
\begin{enumerate}
  \item Save $\textit{result} \gets \textit{heap}[0]$.
  \item $\textit{heap}[0] \gets \textit{heap}[n-1]$; decrement $n$.
  \item $i \gets 0$. While $i$ has a child smaller than itself: let
        $\textit{smallerChild}$ be the smaller of its two children (only
        those that exist); swap $\textit{heap}[i]$ with
        $\textit{heap}[\textit{smallerChild}]$; $i \gets
        \textit{smallerChild}$.
  \item Return \textit{result}.
\end{enumerate}

\begin{algorithm}[H]
\caption{Min Heap: Push and Pop}
\begin{algorithmic}[1]
\Procedure{Push}{$x$}
  \State append $x$; $i \gets n - 1$
  \While{$i > 0$ \textbf{and} $\textit{heap}[(i-1)/2] > \textit{heap}[i]$}
    \State swap $\textit{heap}[(i-1)/2]$ and $\textit{heap}[i]$
    \State $i \gets (i-1)/2$
  \EndWhile
\EndProcedure
\Procedure{Pop}{}
  \State $\textit{result} \gets \textit{heap}[0]$
  \State $\textit{heap}[0] \gets \textit{heap}[n-1]$; remove last slot
  \State $i \gets 0$
  \While{$i$ has a smaller child}
    \State $j \gets$ index of the smaller existing child of $i$
    \If{$\textit{heap}[i] \le \textit{heap}[j]$} \textbf{break} \EndIf
    \State swap $\textit{heap}[i]$ and $\textit{heap}[j]$; $i \gets j$
  \EndWhile
  \State \Return \textit{result}
\EndProcedure
\end{algorithmic}
\end{algorithm}

\subsection*{Dry run}

\dryrun{Push $5, 3, 8, 1$ in order, then pop once.}

\begin{itemize}
  \item Push $5$: heap $=[5]$.
  \item Push $3$: append, heap $=[5,3]$. Compare $3$ to parent $5$:
        $3<5$, swap. Heap $=[3,5]$.
  \item Push $8$: append, heap $=[3,5,8]$. Parent of index $2$ is index
        $0$, value $3$. $8 \ge 3$, no swap. Heap $=[3,5,8]$.
  \item Push $1$: append, heap $=[3,5,8,1]$. Parent of index $3$ is index
        $1$, value $5$. $1<5$, swap: heap $=[3,1,8,5]$. New index $1$;
        parent is index $0$, value $3$. $1<3$, swap: heap $=[1,3,8,5]$.
  \item Pop: result $=1$. Move last element ($5$) to root: heap
        $=[5,3,8]$. At index $0$ ($5$): children at indices $1,2$ are
        $3,8$; smaller is $3$ at index $1$. $5>3$, swap: heap
        $=[3,5,8]$. At index $1$ ($5$): only possible child would be at
        index $3$, out of range. Stop.
\end{itemize}

The pop returns $1$ (correctly the minimum), and the heap array
$[3,5,8]$ still satisfies the min-heap invariant.

\subsection*{C++ implementation}

\begin{lstlisting}
#include <bits/stdc++.h>
using namespace std;

class MinHeap {
    vector<int> heap;

    void siftUp(int i) {
        while (i > 0 && heap[(i - 1) / 2] > heap[i]) {
            swap(heap[(i - 1) / 2], heap[i]);
            i = (i - 1) / 2;
        }
    }

    void siftDown(int i) {
        int n = heap.size();
        while (true) {
            int left = 2 * i + 1, right = 2 * i + 2, smallest = i;
            if (left < n && heap[left] < heap[smallest]) smallest = left;
            if (right < n && heap[right] < heap[smallest]) smallest = right;
            if (smallest == i) break;
            swap(heap[i], heap[smallest]);
            i = smallest;
        }
    }

public:
    void push(int x) {
        heap.push_back(x);
        siftUp(heap.size() - 1);
    }

    int pop() {
        int result = heap[0];
        heap[0] = heap.back();
        heap.pop_back();
        if (!heap.empty()) siftDown(0);
        return result;
    }

    int top() { return heap[0]; }
    bool empty() { return heap.empty(); }
};

int main() {
    MinHeap h;
    h.push(5); h.push(3); h.push(8); h.push(1);
    cout << "Popped: " << h.pop() << endl;
    cout << "New top: " << h.top() << endl;
    return 0;
}
\end{lstlisting}

\begin{complexitybox}
\textbf{Time Complexity.} Both \textit{push} and \textit{pop} touch only
one root-to-leaf path, of length $O(\log n)$, doing $O(1)$ work per level,
giving
\[
  O(\log n) \text{ per push/pop}, \qquad O(1) \text{ per top}.
\]
\textbf{Space Complexity.} $O(n)$ to store the $n$ elements; no extra
auxiliary structures are needed beyond the array itself.
\end{complexitybox}

\begin{takeawaybox}
Every problem in the rest of this chapter uses \texttt{std::priority\_queue}
directly rather than reimplementing this class, but understanding sift-up
and sift-down by hand here is what makes the complexity claims in every
later section ($O(\log n)$ per heap operation) something you can justify
rather than merely recite.
\end{takeawaybox}
